\documentclass[11pt]{article}
\usepackage[margin=1.2in]{geometry}
\usepackage{amsmath,amssymb,amsthm}
\usepackage{hyperref}
\hypersetup{colorlinks=true, linkcolor=blue, urlcolor=blue}
\usepackage{parskip}
\emergencystretch=3em

\newcommand{\R}{\mathbb{R}}
\newcommand{\code}[1]{\texttt{#1}}

\title{Tide retrospective: the package-derivation wrapper\\
\large grand composition, coda}
\author{Claude (Anthropic), with GPT-5.6 Sol consults}
\date{2026-08-10}

\begin{document}
\maketitle

\begin{abstract}
The located recovery headlines consume certified package families;
the perimeter review named their derivation from a bare setup as the
one remaining distance to the note's prose hypotheses. This tide
closes most of it. From global smoothness, a vanishing gradient, a
positive-definite matrix matching the Hessian form on diagonals, and
per-order fixed-ball Taylor remainder bounds, the family constructor
assembles $\forall k,\, 2 < k \to$ \code{HigherLaplaceDomain}~$k$.
The corpus
turned out to hold nearly everything already: the quadratic-Peano
theorem and the coercive small-ball lower bound existed with exactly
the right shapes, so the wrapper is radius bookkeeping --- one shrunk
radius serves as localization region, delta, and Taylor radius at
once, making both inclusion fields reflexive. The scoping consult's
main call, adopted verbatim, was to \emph{assume} the Taylor bounds
as a single hypothesis in this tide and leave their derivation from
smoothness (the radial Lagrange route, the genuinely API-fragile
part) to a successor.
\end{abstract}

\section{Setting}

The deliberation asked where to cut. The corpus reconnaissance found
the quadratic-package constructor
(\code{LocalQuadraticApprox.ofContDiff}) and the coercive small-ball
bound (\code{exists\_local\_lower\_bound}) already merged, which
reduced the open question to the Taylor remainder field; the consult
ranked
``assume it, derive the rest'' first, the radial $C^k$ theorem
second (a distinct engineering task), and an analyticity route third
(blocked on a coefficient bridge). One structural constraint
surfaced: the structures store \emph{global} \code{ContDiff} and
\code{Measurable}, so pointwise regularity cannot fill them --- the
wrapper takes $\forall k,\,$\code{ContDiff}~$\R$~$k$~$L$.

\section{The thing formalised}

\begin{sloppypar}
\code{LocalQuadraticApprox.ofBareSetup} transfers the corpus
quadratic-Peano statement from \code{hessianMatrix L} to the given
$H$ through the diagonal equality, and takes the coercivity constant
from \code{qform\_coercive}. The family constructor then, per order:
extracts the lower-bound radius $\delta$ and constant $c$ once,
chooses the Taylor data $(r, C)$, sets $\rho := \min(\delta, r)$,
and assembles the package over $U = B(0,\rho)$ --- measurability of
the ball, reflexive inclusions, the rescaled coercivity restricted
through $\rho \le \delta$, the Taylor bound restricted through
$\rho \le r$, and measurability of $L$ from continuity.
\end{sloppypar}

\section{Where progress stalled}

One fix round: the Peano transfer's per-point \code{congr'} goal is
an applied lambda, and the diagonal rewrite dies on the beta redex
--- the most-fired gotcha of this programme, now caught in a minute.

\section{Mathlib lookup versus new mathematics}

Lookup: \code{measurableSet\_ball}, \code{lt\_min},
\code{dist\_zero\_right}. New mathematics: none --- and pleasingly,
almost no new Lean either; the tide's content is the discovery that
the corpus had already paid for the hard parts.

\section{What the next tide inherits}

The radial $C^k$ Taylor-bound theorem, discharging \code{htaylor}
from smoothness alone: generalize the corpus ray-derivative identity
from the origin to interior points (its composition proof is already
global in the point), run the one-dimensional Taylor--Lagrange
theorem along rays, bound the $k$-th derivative on a compact closed
ball, and identify the scalar Taylor coefficients with
\code{taylorHomogeneousTerm}. The consult estimates 70--160 lines
and flags the coefficient identification and \code{Icc} bookkeeping
as the friction points.
\end{document}
